% !TEX TS-program = xelatex
% !TEX encoding = UTF-8 Unicode
% !Mode:: "TeX:UTF-8"

\documentclass{resume}
\usepackage{zh_CN-Adobefonts_external} % Simplified Chinese Support using external fonts (./fonts/zh_CN-Adobe/)
% \usepackage{NotoSansSC_external}
% \usepackage{NotoSerifCJKsc_external}
% \usepackage{zh_CN-Adobefonts_internal} % Simplified Chinese Support using system fonts
\usepackage{linespacing_fix} % disable extra space before next section
\usepackage{cite}
\usepackage{xcolor} % for custom colors
\usepackage{framed} % for shaded boxes
\usepackage{mdframed} % for advanced frames

% 定义灰色底纹环境
\definecolor{lightgray}{gray}{0.95}
\newenvironment{graybox}{%
  \begin{mdframed}[backgroundcolor=lightgray,linewidth=0pt,innerleftmargin=8pt,innerrightmargin=8pt,innertopmargin=6pt,innerbottommargin=6pt]
}{%
  \end{mdframed}
}

\begin{document}
\pagenumbering{gobble} % suppress displaying page number

\name{王岗}

\basicInfo{
  \email{acwg345@163.com} \textperiodcentered\ 
  \phone{(+86) 186-5188-5169} \textperiodcentered\
  % \textbf 党员 \textperiodcentered\
  \github[VincentWong1]{https://github.com/VincentWong1}}
 
\section{\faGraduationCap\  教育背景}
\datedsubsection{\textbf{南京⼤学\ ⼯程管理学院}}{2013 -- 2016}
\textit{工学硕士}\ ⼯业⼯程，\ GPA: 4.22/5.00
\datedsubsection{\textbf{南京大学\ 电⼦科学与⼯程学院}}{2009 -- 2013}
\textit{工学学士}\ 通信工程，\ GPA: 4.09/5.00

\small{\datedline{\textit{辅修}\ 商学院经济学院 - 金融学}{2010 -- 2011}}

\section{\faUsers\ 工作经历}
\datedsubsection{\textbf{阿里巴巴}\ 杭州}{2021年3月 -- 至今}

\begin{graybox}
\datedline{\role{中国电商事业群-淘宝闪购-技术中心-供给技术部}{高级应用算法专家} {\small\textcolor{gray}{虚线带算法团队12人}}}{2023年8月 -- 至今}

\vspace{0.5em}
\noindent{\color{gray!40}\rule{\textwidth}{0.2pt}}
\textbf{Agent-0：在闪购商家营销活动配置场景的autoHarness实践}\ {\small\textcolor{gray}{\hfill 202502-至今,\ 项目负责人}}

\vspace{0.3em}
{\small\textbf{项目背景与定位}}
\begin{onehalfspacing}
\begin{itemize}
  \item 主导“商家营销活动推荐与配置”场景的端落地，主牵 Harness-Model 一体化微调方案的落地。
\end{itemize}
\end{onehalfspacing}

\vspace{0.3em}
{\small\textbf{技术架构设计}}
\begin{onehalfspacing}
\begin{itemize}
  \item \textbf{Harness Engineering 层}：设计 DeepAgents + DSPy 混合架构，以 DeepAgents 开源框架完成单场景 Agent 原型的快速搭建与验证，随后接入 DSPy 框架引入 Self-Verification 驱动机制，实现 Harness 节点的自优化闭环迭代——以数据驱动的自优化范式取代传统依赖人工经验的手工方式。
  \item \textbf{Domain-LLM 训练评测层}：设计 DeepEval + Langfuse 组合评测体系，构建 Agent Runtime 整体测评流水线——DeepEval 负责自动化数据跑测，Langfuse 负责运行轨迹的可视化分析与归因，实现对模型能力边界与 Harness 设计效果边界的解耦与定位。
\end{itemize}
\end{onehalfspacing}

\vspace{0.5em}
\noindent{\color{gray!40}\rule{\textwidth}{0.2pt}}
\textbf{EBM饿了么B端垂域通用大模型研发}\ {\small\textcolor{gray}{\hfill 202405-202501,\ 负责人}}

\vspace{0.3em}
{\small\textbf{项目背景与定位}}
\begin{onehalfspacing}
\begin{itemize}
  \item 饿了么B端垂域解空间的有效搜索依赖对垂域数据分布的有效拟合，基于此探索qwen系列小尺寸模型在垂域应用上的有效落地路径。
\end{itemize}
\end{onehalfspacing}

\vspace{0.3em}
{\small\textbf{核心贡献}}
\begin{onehalfspacing}
\begin{itemize}
  \item \textbf{主导了 EBM 系列模型全周期研发与规模化落地}：主导了 EBM-R1（垂域推理）、EBM-R1-MoE（混合架构推理）、EBM-R1-$\pi$（内化 Python 解题的 PAL 模型）、EBM-Instruct-MoE（蒸馏+垂域偏好对齐）四个模型版本的研发。线上覆盖 B 端 85\% token 量（约 420 亿 token），较调用百炼 DeepSeek-R1 日均节约成本 2.7 万元；垂域知识 +7.97\%、推理 +11.60\%、指令 +3.33\%，而通用能力仅推理指标出现有限衰退 -3.32\%;think-model 平均 token 长度较基座下降 27\%，以有限的效果成本换取了大幅的效率提升。
  \item \textbf{建立 EBM 训练与发布规范}：构建 Pre-training（CPT + MID）与 Post-training（SFT + RFT + RL）Pipeline，建立“离线数据生产 $\rightarrow$ 在线大规模并行训练（MDL）$\rightarrow$ 规模化测评（OpenLM）”标准研发范式。
  \item \textbf{主导大模型全栈基建开发}：设计了 Adversarial Dual-Agent 做垂域复杂推理数据合成；建设 E-LLM-Eval 测评体系，采用“垂域/通用 $\times$ 知识/推理/指令 $\times$ 餐饮/零售/其他”多级多维结构，兼顾通用与垂域表现。
  \item \textbf{训练过程优化}：SFT 阶段基于 MoDS 方法原型，用 CrossEntropy 排序代替原 Greedy Selection，实现线上应用数据的自动化过滤回流（较随机采样，垂域推理 +2.55pt、指令 +1.79pt、安全 +2.50pt）；RL 阶段设计了“在/离线拒绝采样+难题回流”机制，提升了学习效率的同时保障了难题探索能力。
\end{itemize}
\end{onehalfspacing}

\vspace{0.3em}
\noindent{\color{gray!40}\rule{\textwidth}{0.2pt}}
\textbf{EBM-AIDA: 饿了么B端自动化数据洞察智能体}\ {\small\textcolor{gray}{\hfill 202405-202412,\ 技术方案设计}}
\begin{onehalfspacing}
\begin{itemize}
  \item 负责了 Trajectory 过程奖励实现的方案设计和理论验证：设计了店铺经营诊断场景下的“元思考范式”，定义了典型 insight 类型（例如 $\langle$DATA\_COMFIRMATION$\rangle$、$\langle$HYPOTHESIS\_TESTING$\rangle$ etc.）。
  \item 将原本无度扩展的 long-running steps 抽象为 FSM（HMM），将每轮推理维持在较稳定的 context span；基于 Pareto 原则仅提取每轮主因，提升解空间的搜索效率；设计了即时奖励 和 累计奖励的混合模式来平衡模型对短期和长期动作的偏好。
  \item 累计25步时，AIDA 相比 ReAct，有效子结论提升8.42，数值幻觉下降24.13，指标下钻深度提升2.67。
\end{itemize}
\end{onehalfspacing}

% \vspace{0.5em}
\noindent{\color{gray!40}\rule{\textwidth}{0.2pt}}
\textbf{饿了么商家助手}\ {\small\textcolor{gray}{\hfill 202308-202404,\ 算法负责人}}

\vspace{0.3em}
{\small\textbf{项目背景与定位}}
\begin{onehalfspacing}
\begin{itemize}
  \item 与工程团队联合孵化了B端首个智能体产品，在行业内较早的探索了以 Agent 助理替代 Chatbot 的可行性。至今已覆盖零售、品牌、城代、官代等业务线 50+ 场景，月活商家 120 万，月均提问量 720 万。
\end{itemize}
\end{onehalfspacing}

\vspace{0.3em}
{\small\textbf{核心贡献}}
\begin{onehalfspacing}
\begin{itemize}
  \item \textbf{配合工程完成了 E-Agent 架构四次迭代}：全程参与了架构演进设计，逐步实现了 agent 架构工作重心从 Workflow 设计转向 LLM 护栏定义：RAG $\rightarrow$ XAgent 多级嵌套 $\rightarrow$ ReAct 快慢并行 $\rightarrow$ Skill 外包，并主导了相关算法基建工作（例如，意图分类、多模态知识库、检索机制、长短期记忆、全链路推荐……），并和工程同学一起沉淀出 E-Agent 行业解决方案，支撑了业务的规模化扩展。
  \item \textbf{设计并主导 IR Agent 开发}：将 Rewrite + RAG 模式升级为 RAG-Fusion 模式，采用 Dense + Sparse + Multi-vector 多路召回，区分 QA-Chunk 和 Doc-Chunk 双路并行，经 RRF 统一排序后，从 chunk 还原 doc 再取Top 6注入 LLM 完成答案生成；测评集召回有效率从 68.35\% 提升至 89.43\%。
  \item \textbf{设计并开发了全链路推荐问系统}：用“知识库原始 Query”和“BERTopic Pipeline 聚合的线上问题”构建动态推荐问池，采用了多窗口（30/7/1d）的流行度权重，适配事前 / 中 / 后的差异化推荐需求。
\end{itemize}
\end{onehalfspacing}
\end{graybox}

\vspace{0.3em}
% \noindent\rule{\textwidth}{0.4pt}
% \vspace{0.3em}
\begin{graybox}
\datedline{\role{淘天集团-客户满意中心-智能研发-人机协同}{高级应用算法专家} {\small\textcolor{gray}{带XP工作台小组3-5人}}}{2021年3月 -- 2023年7月}

\vspace{0.5em}
\noindent{\color{gray!40}\rule{\textwidth}{0.2pt}}

\textbf{智能XP工作台专项}\ {\small\textcolor{gray}{\hfill 202202-202307,\ 算法负责人}}

\vspace{0.3em}
{\small\textbf{基于LLM的工作台算法能力升级探索}}
\begin{onehalfspacing}
\begin{itemize}
  \item 和手淘团队合作共建通用大模型，负责提供客服领域标注数据；做了 LoRA 微调版本，在小蜜 FAQ 问答、自营小蜜商品售前问答、客满技术质检等场景落地应用。
  \item 在2w条域内客服域数据上对ChatGLM-6B做 CT，实现了模型基于“context+会员进线状态”去理解问题、回复特定话术的定制效果。
\end{itemize}
\end{onehalfspacing}

\vspace{0.3em}
{\small\textbf{智能质检}}
\begin{onehalfspacing}
\begin{itemize}
  \item 整合“服务一致性”质检项，将“服务态度、问题定位、问题解决”原离线服务改造为事中的实时算法服务，实现对客服小二服务过程的及时修正。
  \item 主推“智能升级”和“转离完结自动化”两套辅助能力：其中“智能升级”通过算法识别高风险场景，自动填报升级工单；“转离完结自动化”通过算法识别克客服承诺并创建待办，系统监督履约。
  \item “智能升级”项目针对0/1分布失衡的小样本学习场景，采用 RoBERTa Distiller+soft-label 混合模式，单场景仅需200-500条标注样本；“转离完结自动化”部分通过 ELECTRA fine-tune+DualCL 增强承诺句特征表达差异，提升边际case的识别准确性（Precision 0.93）。
\end{itemize}
\end{onehalfspacing}

\vspace{0.3em}
\textbf{主动服务专项}\ {\small\textcolor{gray}{\hfill 202103-202201,\ 算法负责人}}
\begin{onehalfspacing}
\begin{itemize}
  \item 建立分人群、分场域的服务卡点识别能力；通过服务权益的uplift预测和流失预警，匹配补偿动作；最终通过\textbf{智能触达机器人（Agent）}实现有心智的主动服务触达。
  \item 承担智能触达机器人开发，基于 FSM 实现了特定业务场景下的智能外呼对话（日均外呼200，会话完成率64\%）；算法层面，基于 BERT fine-tune 实现意图分类；在 FSM 对话流基础上集成 QQ-Bot 模块，采用 DSSM 模型完成 Query-Question 匹配，ACC@1 达 0.87，补充领域 FAQ 的精准响应能力。
  \item 从线上语聊挖掘高频对话序列，降低对话流的生产成本，为 TSBC 架构下无流程指引的自有对话实现提供数据支撑。
\end{itemize}
\end{onehalfspacing}
\end{graybox}

\vspace{0.5em}
% \noindent\rule{\textwidth}{0.4pt}
\datedsubsection{\textbf{苏宁金融}\ 南京}{2016年7月 -- 2021年2月}

\begin{graybox}
\datedline{\role{苏宁金融科技事业部-智能营销产品中心-算法研发部}{经理} {\small\textcolor{gray}{管理南京10人，上海3人}}}{2016年7月 -- 2021年2月}

\vspace{0.5em}
\noindent{\color{gray!40}\rule{\textwidth}{0.2pt}}

\textbf{智能营销产品中心}\ {\small\textcolor{gray}{\hfill 201901-202102,\ 经理}}
\begin{onehalfspacing}
\begin{itemize}
  \item \textbf{CRM的实时行业画像功能}：Randomized Lasso 做特征筛选，Spark Streaming 实现用户画像实时分析。
  
  \item \textbf{首页理财活动CTR预测}：利用 PCA 降维+稠密化处理实现多渠道特征融合，GBDT+LR 做点击预测。
  
  \item \textbf{营销分}：LightGBM 做模型训练，Platt Scaling 做多模型融合，实现评分模型的客群全覆盖。
\end{itemize}
\end{onehalfspacing}

\vspace{0.3em}
\textbf{智能客服项目组}\ {\small\textcolor{gray}{\hfill 201901-202006,\ 算法专家}}
\begin{onehalfspacing}
\begin{itemize}
  \item \textbf{客服机器人开发}：负责框架设计、意图识别（fastText，macro precision 87.5\%）、相似问匹配（WMD+Levenshtein）、镜像商业化输出
  
  \item \textbf{催收机器人开发}：设计“FSM 问答转移 + slotfilling 动态加载”机制，接入 TTS 实现智能语音催收。
  
  \item \textbf{质检机器人开发}：离线训练 Character-level TextCNN 做高危词语境含义识别（多义高危词 F1 0.70）。
\end{itemize}
\end{onehalfspacing}

\vspace{0.3em}
\textbf{智能投顾项目组}\ {\small\textcolor{gray}{\hfill 201607-201812,\ 算法高级工程师}}
\begin{onehalfspacing}
\begin{itemize}
  \item \textbf{“苏宁智投”产品的底层算法开发}：合作开发"苏宁智投"产品，负责基金风险收益计算、基金池筛选及资产归类模块开发；协作开发基于遗传算法的均值-方差最优化计算模块（年化收益率 7.19-19.84\%）。
  
  \item \textbf{智能营销产品推荐算法开发}：特定业务场景下的 LR+评分卡营销模型开发，共同拟定了"基于逻辑回归的评分卡建模"标准流程。
\end{itemize}
\end{onehalfspacing}
\end{graybox}

% Reference Test
%\datedsubsection{\textbf{Paper Title\cite{zaharia2012resilient}}}{May. 2015}
%An xxx optimized for xxx\cite{verma2015large}
%\begin{itemize}
%  \item main contribution
%\end{itemize}

\vspace{0.3em}
\section{\faCogs\ IT 技能}
% increase linespacing [parsep=0.5ex]
\begin{itemize}[parsep=0.5ex]
  \item 编程语言: Python == SQL > Shell > R
  \item 开发: PyTorch, TensorFlow, Spark, Hadoop, Hive, Redis, ES, HBase, MySQL
  \item AI工具/框架：Claude Code, OpenClaw, DeepAgents, DSPy, Langfuse/LangSmith, DeepEval, XAgent etc.
  \item 训练框架：TuningFactory(/LlamaFactory), ROLL
  \item 算法: 强化学习（PPO/GRPO/DAPO), 监督微调（Full/LoRA)，深度学习（CNN/RNN/transformer)，机器学习
\end{itemize}

\section{\faHeartO\ 作品}
\textbf{论文}
\begin{itemize}[parsep=0.3ex]
  \item \textit{} Towards Autonomous Business Intelligence via Data-to-Insight Discovery Agent. ICML under review. \hfill 2026
\end{itemize}
\textbf{专利}
\begin{itemize}[parsep=0.3ex]
  \item \textbf{CXTA177467} - 一种数据样本的生成方法、问答方法和电子设备 \hfill 12/2025 \\
  {\small\textcolor{gray}{（基于多智能体对抗的垂域数据合成）}}
  
  \item \textbf{CXTA177008} - 问答数据生成方法、问答方法、计算机设备和存储介质 \hfill 12/2025 \\
  {\small\textcolor{gray}{（基于跨文档知识图谱的领域推理问生产）}}
  
  \item \textbf{CXTA151576} - 任务型对话方法、存储介质及计算机设备 \hfill 06/2024 \\
  {\small\textcolor{gray}{（E-Agent：兼顾速度、准确性、稳定性、和易于改进的智能助手方案）}}
  
  \item \textbf{CXTA115006} - 待办任务创建、履约动作分类模型构建方法、电子设备 \hfill 03/2023 \\
  {\small\textcolor{gray}{（客服领域中基于自然语言识别的承诺待办系统及其算法和流程）}}

  \item \textbf{CN114374770A} - ⼀种语⾳质检⽅法、装置、计算机设备及存储介质 \hfill 04/2021
  \item \textbf{CN114333813A} - 可配置智能语⾳机器⼈的实现⽅法、装置和存储介质 \hfill 04/2021
  \item \textbf{CN111898377A} - ⼀种情感识别⽅法、装置、计算机设备及存储介质 \hfill 07/2020
  \item \textbf{CN112036906A} - ⼀种数据处理⽅法、装置、设备 \hfill 07/2020
  \item \textbf{CN110298759A} - ⼀种基⾦诊断⽅法、装置及计算机可读存储介质 \hfill 2019
\end{itemize}

\section{\faInfo\ 其他}
% increase linespacing [parsep=0.5ex]
\begin{itemize}[parsep=0.5ex]
  \item 语言: IELTS (6.5), TOEFL (93), GRE (Quantitative:170), 中级⼝译, CET6
  \item 管理: PMP 证书
\end{itemize}

%% Reference
%\newpage
%\bibliographystyle{IEEETran}
%\bibliography{mycite}
\end{document}
